\underline{Demostración.}

\begin{enumerate}[label=(\arabic*)]
\item Si
\[
 0\longrightarrow M'\longrightarrow M\longrightarrow M''
 \longrightarrow0
\]
es una sucesión exacta de $A$-módulos, entonces, para toda álgebra de
polinomios $P$ sobre $R$, la sucesión
\[
 0\longrightarrow\operatorname{Der}_R(P,M')
 \longrightarrow\operatorname{Der}_R(P,M)
 \longrightarrow\operatorname{Der}_R(P,M'')
 \longrightarrow0
\]
es exacta. Por tanto, según 3.3(c),
\[
 0\longrightarrow\operatorname{Der}_R(-,M')
 \longrightarrow\operatorname{Der}_R(-,M)
 \longrightarrow\operatorname{Der}_R(-,M'')
 \longrightarrow0
\]
es una sucesión exacta de haces.

\item El funtor $\operatorname{Der}_S(-,M)$ es un haz aditivo sobre
$\underline C_{A|S}$, y
\[
 \operatorname{Der}_R(B,M)
 =\ker\bigl(\operatorname{Der}_S(B,M)
 \longrightarrow\operatorname{Der}_S(R,M)\bigr)
\]
para todo $B\in\operatorname{Ob}(\underline C_{A|R})$.
La afirmación se deduce entonces de 3.7(c).

\item Sea
\[
 S:\underline C_{A|R}\longrightarrow
 \underline C_{A\otimes_RS|S},
 \qquad X\longmapsto X\otimes_RS.
\]
Entonces
$\widetilde S\operatorname{Der}_S(-,M)
=\operatorname{Der}_R(-,M)$, de modo que la afirmación se deduce de
3.7(a).

%%%%%%%%%%  scan idx 99  |  folio 88  %%%%%%%%%%

\item Esto se deduce de 3.9. En efecto,
$F=\operatorname{Der}_R(-,N)$ es un haz aditivo sobre
$\underline C_{S^{-1}A|R}$, y
$F_A\in\operatorname{Ob}
(\widetilde{\underline C}_{S^{-1}A|A})$ también es aditivo, pues
\begin{align*}
 F_A(X)
 &=\ker\bigl(F(X)\longrightarrow F(A)\bigr)\\
 &=\ker\bigl(\operatorname{Der}_R(X,N)
   \longrightarrow\operatorname{Der}_R(A,N)\bigr)\\
 &=\operatorname{Der}_A(X,N).
\end{align*}
Así, $F_A=\operatorname{Der}_A(-,N)$.

\item La primera afirmación se deduce de 3.8(a). Supongamos ahora que
$\operatorname{Tor}_i^R(A,A)=0$ para todo $i>0$ y que
\[
 \Delta:\operatorname{Spec}(A)\longrightarrow
 \operatorname{Spec}(A)\times_{\operatorname{Spec}(R)}
 \operatorname{Spec}(A)
\]
es una inmersión abierta. Sea
$\varphi:A\otimes_RA\longrightarrow A$ el morfismo de multiplicación.
Para todo ideal maximal $\mathfrak m$ de $A$, el morfismo inducido
\[
 (A\otimes_RA)_{\varphi^{-1}(\mathfrak m)}
 \longrightarrow A_{\mathfrak m}
\]
es un isomorfismo. La primera afirmación, junto con 3.12(4), implica
entonces que el morfismo diagonal
\[
 \Delta_{\mathfrak m}:D^i(A|R,M)_{\mathfrak m}
 \longrightarrow D^i(A|R,M)_{\mathfrak m}
 \oplus D^i(A|R,M)_{\mathfrak m}
\]
es un isomorfismo para todo ideal maximal $\mathfrak m$ de $A$.
Por consiguiente,
\[
 \Delta:D^i(A|R,M)\longrightarrow
 D^i(A|R,M)\oplus D^i(A|R,M)
\]
es un isomorfismo, y por tanto $D^i(A|R,M)=0$ para todo $i>0$.

\item En vista de 3.12(4), podemos suponer que $A$ es una
$R$-álgebra de tipo finito. Consideremos el morfismo de haces
\[
 E\otimes_A\operatorname{Der}_R(-,M)
 \longrightarrow
 \operatorname{Der}_R(-,E\otimes_AM)
\]
sobre $\widetilde{\underline C}_{A|R}$. Como $E$ es plano sobre $A$,
el morfismo
\[
 E\otimes_A\operatorname{Der}_R(B,M)
 \xrightarrow{\ \sim\ }
 \operatorname{Der}_R(B,E\otimes_AM)
\]
es un isomorfismo siempre que
$B\in\operatorname{Ob}(\underline C_{A|R})$ sea de tipo finito sobre
$R$. Según 3.3(b), se sigue que
\[
 D^i\bigl(A|R,E\otimes_A\operatorname{Der}_R(-,M)\bigr)
 \longrightarrow
 D^i\bigl(A|R,\operatorname{Der}_R(-,E\otimes_AM)\bigr)
\]
es un isomorfismo. Por otra parte, $E\otimes_A-$ es exacto y, por tanto,
\[
 E\otimes_AD^i(A|R,M)
 =D^i\bigl(A|R,E\otimes_A\operatorname{Der}_R(-,M)\bigr).
\]
Esto demuestra la afirmación.

\item Esto se deduce de 3.3(a).

%%%%%%%%%%  scan idx 100  |  folio 89  %%%%%%%%%%

\item En vista de 3.10 y 3.11, basta mostrar que
\[
 \operatorname{Coker}\bigl(
 \operatorname{Der}_R(P,M)\longrightarrow
 \operatorname{Hom}_A(I/I^2,M)\bigr)
 =\operatorname{Exalcomm}(A|R,M),
\]
donde $\operatorname{Exalcomm}(A|R,M)$ denota el conjunto de clases de
isomorfismo de extensiones de $R$-álgebras de $A$ por $M$. Sea
\[
 0\longrightarrow M\longrightarrow A'\longrightarrow A
 \longrightarrow0
\]
una extensión de este tipo. Como $P$ es un álgebra de polinomios sobre
$R$, obtenemos un diagrama conmutativo
\[
\begin{tikzcd}[column sep=large]
0 \arrow[r] & I \arrow[r] \arrow[d]
& P \arrow[r] \arrow[d]
& A \arrow[r] \arrow[d,equal] & 0 \\
0 \arrow[r] & M \arrow[r]
& A' \arrow[r]
& A \arrow[r] & 0.
\end{tikzcd}
\]
El morfismo $I\longrightarrow M$ está determinado por la extensión
$A'$, salvo la acción de $\operatorname{Der}_R(P,M)$. Obtenemos así un
morfismo
\[
 \Phi:\operatorname{Exalcomm}(A|R,M)
 \longrightarrow
 \operatorname{Coker}\bigl(
 \operatorname{Der}_R(P,M)\longrightarrow
 \operatorname{Hom}_A(I/I^2,M)\bigr).
\]
Recíprocamente, un morfismo $P$-lineal $I\longrightarrow M$ produce,
mediante un pushout, una extensión que forma parte de un diagrama
\[
\begin{tikzcd}[column sep=large]
0 \arrow[r] & I \arrow[r] \arrow[d]
& P \arrow[r] \arrow[d]
& A \arrow[r] \arrow[d,equal] & 0 \\
0 \arrow[r] & M \arrow[r]
& A' \arrow[r]
& A \arrow[r] & 0.
\end{tikzcd}
\]
Dos morfismos $P$-lineales $I\rightrightarrows M$ que difieren en un
elemento de $\operatorname{Der}_R(P,M)$ producen extensiones isomorfas.
Por tanto, hay un morfismo
\[
 \Psi:\operatorname{Coker}\bigl(
 \operatorname{Der}_R(P,M)\longrightarrow
 \operatorname{Hom}_A(I/I^2,M)\bigr)
 \longrightarrow\operatorname{Exalcomm}(A|R,M).
\]
Se comprueba inmediatamente que $\Phi$ y $\Psi$ son inversos uno de
otro.
\end{enumerate}
\hfill$\square$
